\section{Ergebnisse}

\subsection{Repräsentative End-to-End-Artefaktkette}

Der vollständige Funktionsnachweis wird anhand eines intern konsistenten
Referenzlaufs dargestellt. Die endgültige Fassung zeigt dafür dieselbe
Verarbeitungskette in aufeinander abgestimmten Ansichten: Rohframe,
\texttt{default}/\texttt{middle}-Maske, ideale Bild-/Maskendomäne,
COLMAP-Sparse-Modell, objektbezogene STS-Gaussians, Route-A-Coarse-Mesh sowie
lokale Centerline und Grad-10-B-Spline.

\begin{quote}
\textbf{Offener Abbildungsnachweis:} Vor der finalen Abgabe sind die
Artefakte eines ausgewählten Golden Runs in identischen beziehungsweise
nachvollziehbar zugeordneten Ansichten zu exportieren. Eine Mischung aus
verschiedenen Matrixläufen darf nicht als einzelne End-to-End-Kette erscheinen.
\end{quote}

Ein archivierter Low-Smoke-Test weist die Pipeline ab dem Eintritt in die
ideale, verzeichnungsfreie Bildszene nach: Ab dieser Stufe sind Bild-, Masken-
und Kamerakoordinaten untereinander konsistent, sodass alle nachgelagerten
Stufen (STS, Meshroute, Rendering, Metriken, Centerline) auf derselben Domäne
arbeiten. Der Lauf verwendet 240 Kameras; davon dienen 210 dem Training und 30
der Evaluation. Die Evaluationsbilder sind gerade dadurch bewertbar, dass sie
nicht am Training beteiligt waren und unabhängig von ihm ausgewertet werden
können. Der Lauf erzeugte Route-A-Mesh, Rendering, objektmaskierte Metriken,
Centerline, B-Spline und GeoJSON. Für den vollständigen Autopilotnachweis ist
zusätzlich ein archivierter Autopilot-Vollauf oder ein gleichwertiger
Vertragstest erforderlich; der Matrixrunner allein belegt nur die
Batchfähigkeit. Die historischen Zwischenstände, die zu den heutigen
Einstellungen in STS, Objektfilterung und SuGaR geführt haben, sind in
Kapitel~\ref{sec:implementierung} (Abschnitt
``Historische Zwischenstände der STS- und Objektfilterungsroute'') sowie in
Abschnitt~\ref{sec:sugar-folgepruefung} dokumentiert.

\subsection{COLMAP-Voruntersuchung}

Ebenso wurde vor den Matrixläufen eine Versuchsreihe mit COLMAP gestartet;
dabei wurde geprüft, inwiefern sich einzelne Einstellungen im Hinblick auf den
zeitlichen Aufwand und die Qualität der Sparse-Rekonstruktion lohnen. Alle
fünf dokumentierten Hauptläufe registrierten 240 von 240 Bildern. Guided
Matching erhöhte die Sparse-Punktzahl gegenüber Plain-SIFT 4096 nur um rund
1,98\,\%, benötigte aber nahezu die doppelte Zeit und erhöhte den mittleren
Reprojektionsfehler. 8192 Merkmale lieferten etwa 0,82\,\% mehr Punkte bei
5,5\,\% mehr Laufzeit. Damit ist Plain-SIFT 4096 ohne Guided Matching für den
Projektstand der begründete Standard. Die vollständige Tabelle steht in
Anhang~\ref{tab:colmap-vorstudie}.

\subsection{Objektmaskierte Bildmetriken}

Ein Matrixlauf ist ein vollständig protokollierter Einzelversuch für genau
eine Kombination aus Kameramodell, FPS, Auflösung und Meshroute. Der Runner
erzeugt dafür einen isolierten Frame- und Maskensatz, führt COLMAP, die ideale
Bilddomäne, STS, die gewählte Meshroute, Rendering und objektmaskierte
Metriken aus und archiviert anschließend die Resultate. In der Entwicklung
traten unter anderem ein durch Unterprozesse konsumierter TSV-Standardinput,
leere feste Eval-Masken sowie ein SuGaR-Import-/Tensorlayoutfehler auf. Diese
Fehler wurden durch stdin-unabhängige Schleifen, einen nichtleeren festen
Eval-Split sowie die Korrektur des Renderhelfers behoben. Dadurch lassen sich
die Matrixläufe heute mit denselben Parametern erneut erzeugen und unter
identischen Kriterien vergleichen.

Die Matrix dient dabei nicht dazu, aus einer einzelnen Bildmetrik eine
„beste“ Pipeline zu küren. Sie zeigt, welche Kombinationen technisch
durchlaufen, welche Fehlerfälle auftreten und wie sich die Ansichtsmetriken
unter kontrollierten Änderungen verhalten. Die Produktionsentscheidung für
Route A stützt sich deshalb vor allem auf die A/B/C/D-Geometrieablation und
den direkten Erhalt der STS-Gaussian-Geometrie; die Matrix ergänzt diesen
Befund um Robustheit und reproduzierbare Vergleichswerte.

Die drei Bildmetriken sind dabei ausschließlich Ansichtsmetriken der jeweils
gerenderten Gaussian-Repräsentation. Die STS-Grafik verwendet
\texttt{sts\_masked.json}; die SuGaR-Grafik verwendet
\texttt{sugar\_coarse\_masked.json}. Ein höherer PSNR- oder SSIM-Wert
beziehungsweise ein niedrigerer LPIPS-Wert bedeutet deshalb zunächst nur eine
ähnlichere Objektansicht zum Ground-Truth-Bild. Daraus folgt weder eine
korrektere Meshoberfläche noch eine bessere Centerline. Für die Meshentscheidung
sind weiterhin Vertex-/Face-Struktur, relative Meshabstände und die visuelle
Artefaktprüfung maßgeblich; eine absolute geometrische Referenz fehlt in der
Projektarbeit.

Ein erneuter vollständiger Matrixlauf ist als Wiederholungsprüfung sinnvoll.
Bei nominal gleichen Parametern sind kleine Abweichungen durch stochastische
Optimierung, Surface-Sampling sowie Software- und Hardwarestände möglich. Eine
Abweichung wird erst dann als technisch relevant bewertet, wenn Commitstände,
Eval-Split, Framezahl, Maskenabdeckung und Seeds übereinstimmen. Für die PA
werden daher nicht alle wiederholten Arbeitskopien als eigenständige
wissenschaftliche Ergebnisse benötigt; entscheidend sind die standardisierten
Metriken, die Endprodukte und die reproduzierbaren Metadaten.

Die zentrale Metrikübersicht ist nun nach Modellstufe getrennt. Die STS-Grafiken
lesen ausschließlich vollständige Original-GS-Route-A-Läufe aus
\texttt{sts\_masked.json}; die SuGaR-Grafiken lesen ausschließlich
\texttt{sugar\_coarse\_masked.json}. Vorhandene STS-Dateien aus später
fehlgeschlagenen SuGaR-Routen bleiben als Fehlernachweis archiviert, werden aber
nicht zusätzlich in der STS-Auswertung dargestellt. Die beiden FPS-Profile und
die jeweilige Runauflösung bleiben sichtbar; ein globales Ranking über
verschiedene Auflösungen wird nicht vorgenommen.

\begin{figure}[H]
  \centering
  \includegraphics[width=\textwidth]{sts_masked_overview.pdf}
  \caption{Getrennte STS-7000-Auswertung der vollständigen Original-GS-Route A
  aus den archivierten \texttt{sts\_masked.json}-Dateien.}
  \label{fig:pa-matrix-overview}
\end{figure}

\begin{figure}[H]
  \centering
  \includegraphics[width=\textwidth]{sts_masked_per_frame_boxplots.pdf}
  \caption[STS-Einzelansichten]{Verteilung der Einzelansichten in den
  vollständigen Original-GS-Route-A-Läufen. Punkte entsprechen den in
  \texttt{per\_frame} archivierten Ansichtsmetriken.}
  \label{fig:pa-sts-boxplots}
\end{figure}

\begin{figure}[H]
  \centering
  \includegraphics[width=\textwidth]{sugar_coarse_masked_overview.pdf}
  \caption{Getrennte SuGaR-Coarse-9000-Auswertung. Die STS-Baseline wird
  nicht als SuGaR-Wert verwendet.}
  \label{fig:pa-sugar-overview}
\end{figure}

Die deutlich niedrigeren SuGaR-Werte gegenüber einigen alten
\texttt{simple\_radial\_sugar}-Zeilen sind kein direkter Widerspruch. Die
alten Zeilen der ursprünglichen Overview bezogen sich trotz ihrer
SuGaR-Variantenbezeichnung auf \texttt{sts\_masked.json}: Der STS-Splat wurde
vor dem fehlgeschlagenen SuGaR-Renderhelfer ausgewertet. Beispielsweise wurden
für den alten 5-FPS-720p-Lauf 29,337 dB PSNR, 0,896 SSIM und 0,122 LPIPS als
STS-Baseline protokolliert. Der entsprechende aktuelle Wert von
\texttt{sugar\_coarse\_masked.json} beträgt 21,156 dB, 0,776 SSIM und 0,152
LPIPS. Dieser Vergleich stellt somit zwei verschiedene Modellstufen
gegenüber: STS-7000 vor SuGaR-Optimierung versus SuGaR-Coarse nach der
zusätzlichen Gaussian-Optimierung und dem SuGaR-Rendering. Die neuen Werte sind
deshalb die fachlich richtigen SuGaR-Coarse-Werte; die alten Werte bleiben als
STS-Baselines historisch nachvollziehbar, dürfen aber nicht als SuGaR-
Ergebnisse bezeichnet werden.

In der jeweiligen Runauflösung zeigt OPENCV-A-Low bei 5 FPS besonders hohe
Bildmetriken. Bei SIMPLE-RADIAL und OPENCV bezieht sich die Auswertung dabei
auf die ideale, gewarpte Bilddomäne; „jeweilige Runauflösung“ bezeichnet nur
die Pixelauflösung innerhalb dieser Domäne. Dieser Befund ist ein
Screening-Ergebnis. Niedrige Auflösung kann Kantenversatz und hochfrequente
Fehler glätten, sodass bessere Bildmetriken nicht automatisch eine bessere
3D-Geometrie bedeuten. Die Werte werden daher nicht als alleinige
Produktionsentscheidung verwendet.

Die Ergebnisse sind nicht monoton. 5 FPS liefern nicht in jeder Kombination
den numerisch besten Wert, und Low-Auflösung ist nicht in jedem Modell besser
als QHD. Die beobachteten Unterschiede entstehen daher wahrscheinlich aus
einer Interaktion von Bildauflösung, zeitlicher Überdeckung, konkreter
Kamerabewegung, Maskenstabilität, Kameramodell und STS-Optimierung. Ein
isolierter Kausalnachweis für einen einzelnen Faktor liegt nicht vor.

\subsection{Produktionskonfiguration aus der Matrixauswertung}
\label{sec:produktionskonfiguration}

Auf Grundlage der Matrixauswertung wurde die Produktionskonfiguration auf
\texttt{OPENCV}, 5 FPS und Route A festgelegt. Bei 720p und 5 FPS erzielen die
\texttt{OPENCV}-A-Läufe die höchsten objektmaskierten STS-Metriken der
untersuchten Kameramodelle (29,61\,dB PSNR, 0,901 SSIM und 0,117 LPIPS
gegenüber 29,37\,dB, 0,896 und 0,123 bei \texttt{SIMPLE\_RADIAL}); die
Konfiguration ist über die archivierten Läufe reproduzierbar belegt. Die
Umstellung bleibt dennoch eine datengestützte Projektentscheidung ohne
unabhängige geometrische Referenz: Bei qhd liegt \texttt{SIMPLE\_RADIAL} mit
28,21\,dB vor \texttt{OPENCV} (27,67\,dB), was die nichtmonotone Interaktion
von Auflösung und Kameramodell unterstreicht. Der Standard gilt daher
ausdrücklich für die Produktionsauflösung 720p bei 5 FPS.

Für diese Produktionskonfiguration ergeben die über die Matrixarchive
gemittelten Laufzeiten je Auflösung rund 2023 Sekunden (720p), 1648 Sekunden
(qhd) und 1141 Sekunden (low). Diese Werte fließen in die Lauf-Preset-Abfrage
des CLI ein (Kapitel~\ref{sec:implementierung}), mit der zwischen
Bildauflösung und Rechenaufwand gewählt werden kann. Als reine Schrittzeiten
von STS bis Postprocessing ersetzen sie keine vollständige End-to-End-
Zeitmessung.

\subsection{Vierfeld-Ablation und Produktionsentscheidung}

Die kontrollierte Ablation trennt Original-STS-Gaussians beziehungsweise
SuGaR-Coarse und zwei Tiefenrouten. Tabelle~\ref{tab:mesh-ablation} enthält die
resultierenden Meshgrößen.

\begin{table}[H]
\centering
\small
\begin{tabularx}{\textwidth}{lXrr}
\toprule
Variante & Eingang und Tiefenroute & Vertices & Faces\\
\midrule
A & Original-STS-GS, Diamond-Mesh-Z-Buffer & 68471 & 135164\\
B & Original-STS-GS, Gaussian-Rasterizer-Tiefe & 63266 & 124976\\
C & SuGaR-Coarse $c=9000$, Diamond-Mesh-Z-Buffer & 79711 & 156935\\
D & SuGaR-Coarse $c=9000$, Gaussian-Rasterizer-Tiefe & 68618 & 135522\\
\bottomrule
\end{tabularx}
\caption{Ergebnisse der Coarse-only-Vierfeld-Ablation bei identischen
Extraktionsparametern und Seed 42.}
\label{tab:mesh-ablation}
\end{table}

A und B liegen großräumig auf derselben Objektgeometrie, sind aber nicht
identisch. Gerichtete Mittelwerte aus jeweils 20000 gesampelten Punkten lagen
bei etwa 1,08 beziehungsweise 0,90\,cm. Gegenüber A betrugen die gerichteten
Mittelwerte für C etwa 1,79 und 3,39\,cm und für D etwa 1,96 und 2,73\,cm. Die
größeren Gegenrichtungsabstände und die visuell stärkere Ausreißerstruktur der
Coarse-Varianten zeigen, dass die zusätzliche Gaussian-Optimierung die
Oberflächenstützung relevant verändert. Diese Werte messen Variantenabstände,
nicht die Genauigkeit zum realen Rohr.

Route A wird deshalb als Produktionsstandard gewählt: Sie bewahrt die
Original-STS-Gaussians als primäre Geometriehypothese und vermeidet eine
zusätzliche Coarse-Optimierung, die für die Centerline-Anwendung mehr Details,
aber auch mehr Spitzen und Randartefakte erzeugen kann. Die Wahl bleibt eine
begründete Projektentscheidung ohne unabhängige Ground Truth. Dass die
Coarse-Route dennoch existiert, erklärt sich aus der SuGaR-Startarchitektur:
Sie ist der dort übliche Weg zur Meshgewinnung, während die direkte Extraktion
ohne Coarse-Optimierung erst im Projektverlauf ergänzend erprobt wurde
(vgl. Kapitel~\ref{sec:grundlagen} sowie Abschnitt~\ref{sec:sugar-folgepruefung}).

\subsection{Abgeschlossene SuGaR-Coarse-Folgeprüfung}
\label{sec:sugar-folgepruefung}

Die zwölf Läufe der separaten Matrix
\texttt{matrix\_sugar\_followup\_12} wurden vollständig ausgeführt. Je sechs
Läufe entfielen auf \texttt{SIMPLE\_RADIAL} und \texttt{OPENCV}; beide
Kameramodelle wurden bei 2 und 5 FPS sowie in 720p, QHD und Low getestet. Alle
zwölf Manifeste tragen den Status \texttt{success}. Damit sind Masken- und
Domänenprüfung, STS-Training, SuGaR-Coarse-Mesh, SuGaR-Coarse-Rendering,
objektmaskierte Metriken und Postprocessing für diese zwölf Konfigurationen
archiviert.

Die Coarse-Route ist der konzeptionelle Standardweg von SuGaR: Meshes werden
dort üblicherweise erst nach einer zusätzlichen Coarse-Optimierung der
Gaussian-Geometrie extrahiert. Die in dieser Arbeit als Produktionsroute
verwendete Route A weicht bewusst von diesem Standardweg ab, weil sie direkt
aus der unveränderten Original-STS-Gaussian-Wolke extrahiert. Diese Variante
wurde erst im Projektverlauf ergänzend erprobt und hat anschließend die
Produktionsrolle übernommen; die Coarse-Route wird deshalb nicht als
Fehlversuch, sondern als wissenschaftlicher Vergleichsarm weitergeführt
(vgl. Kapitel~\ref{sec:grundlagen}). Abbildung~\ref{fig:pa-panel-sugar}
dokumentiert ausgewählte historische Zwischenstände dieser Entwicklung.

\begin{figure}[H]
  \centering
  \includegraphics[width=\textwidth]{pa_panel_sugar_iteration.png}
  \caption{Historische SuGaR-Coarse-Zwischenstände: Meshstufen mit
  Zielzähler 9001 sowie bilaterale Vergleichsläufe im 9000er-/10.000er- und
  10.000er-/15.000er-Bereich. Die Angabe ``9200'' kam ausschließlich in zwei
  historischen Screenshot-Dateinamen vor und ist weder durch Manifest noch
  durch Log bestätigt; die zugehörigen Bilder sind hier einheitlich als
  Stände mit Zielzähler 9001 eingeordnet.}
  \label{fig:pa-panel-sugar}
\end{figure}

Die neuen Grafikquellen trennen die Modellstufen ausdrücklich. Neben den beiden
Übersichten werden die Streuung der Einzelansichten, die gepaarte Differenz
zwischen SuGaR-Coarse und STS sowie die Beziehung zwischen Bildmetrik und
archivierter Laufzeit separat dargestellt. Diese Zusatzgrafiken sind
Diagnose- und Interpretationshilfen, keine neuen Geometriemetriken.

Die Boxplotpunkte liegen jeweils in der x-Kategorie ihres Laufs und werden nur
leicht horizontal verteilt. Die y-Achsen berücksichtigen alle archivierten
Einzelwerte einschließlich Ausreißern. In der Laufzeitgrafik ist die x-Achse
auf 15--40 Minuten begrenzt und zeigt die Gesamtlaufzeit des archivierten
Pipeline-Laufs. Grün und Lila kodieren STS und SuGaR-Coarse; kräftige und
blasse Farbtöne kodieren 5 beziehungsweise 2 FPS. Kreis, Quadrat und Dreieck
stehen für 720p, QHD und Low. Die Kamera wird über die Achsenbeschriftung
angegeben, nicht über eine zusätzliche SR-/OPENCV-Legende. Die separate Matrixstatusgrafik gehört
nicht mehr zum PA-Anhang; der Status bleibt als Tabelle und Manifestinformation
erhalten.

\begin{figure}[H]
  \centering
  \includegraphics[width=\textwidth]{sugar_coarse_masked_per_frame_boxplots.pdf}
  \caption[SuGaR-Coarse-Einzelansichten]{Verteilung der Einzelansichten in
  den zwölf SuGaR-Coarse-Läufen. Punkte entsprechen den in
  \texttt{per\_frame} archivierten Ansichtsmetriken.}
  \label{fig:pa-sugar-boxplots}
\end{figure}

\begin{figure}[H]
  \centering
  \includegraphics[width=\textwidth]{sugar_coarse_vs_sts_delta.pdf}
  \caption{Gepaarte Differenz zwischen SuGaR-Coarse und STS bei identischer
  Follow-up-Konfiguration. Negative PSNR-/SSIM-Werte und positive LPIPS-Werte
  bedeuten in diesen Läufen eine schlechtere SuGaR-Coarse-Bildmetrik; die
  Differenz ist kein Mesh-Genauigkeitsnachweis.}
  \label{fig:pa-sugar-delta}
\end{figure}

\begin{figure}[H]
  \centering
  \includegraphics[width=\textwidth]{metric_vs_runtime.pdf}
  \caption{Explorative Beziehung zwischen objektmaskierten Bildmetriken und
  der Gesamtlaufzeit des archivierten Pipeline-Laufs. SAM3- und COLMAP-Wandzeit
  sind nur enthalten, wenn sie innerhalb des archivierten Laufs lagen.}
  \label{fig:pa-metric-runtime}
\end{figure}

Die Mittelwerte der beiden FPS-Profile liegen für
\texttt{SIMPLE\_RADIAL/SuGaR-Coarse} bei 21,38 dB PSNR, 0,764 SSIM und 0,1534
LPIPS in 720p; bei QHD bei 21,69 dB, 0,761 und 0,1407; bei Low bei 21,89 dB,
0,752 und 0,1261. Für \texttt{OPENCV/SuGaR-Coarse} ergeben sich 21,45 dB,
0,764 und 0,1531 in 720p; 21,62 dB, 0,764 und 0,1367 in QHD sowie 21,81 dB,
0,749 und 0,1255 in Low. Die Werte sind reine Ansichtsmetriken und werden
nicht als Beleg einer besseren 3D-Geometrie interpretiert.

Die \texttt{sugar\_coarse\_masked.json}-Dateien enthalten je nach FPS zwölf
beziehungsweise dreißig Evaluationsansichten. Die zugehörigen
\texttt{sugar\_refined\_masked.json}-Dateien sind in allen zwölf Archiven
wegen eines nicht exportierten \texttt{refined.ply} übersprungen. Der neue
Befund weist daher die technische SuGaR-Coarse-Route nach, schließt aber noch
keinen Vergleich der SuGaR-Refined-Stufe ein.

\subsection{Funktionsnachweis der Endstufen}

Die erfolgreichen A-Läufe erzeugten neben den STS-Metriken auch Coarse-Ply,
kompatibles OBJ, lokale Centerline, B-Spline und GeoJSON. Die Georeferenzierung
war in den Matrixarchiven mangels validierter CloudCompare-Matrix teilweise ein
Translation-Fallback. Damit ist die technische Kette bis zu lokalen
Geometrieprodukten nachgewiesen; eine globale Genauigkeit ist daraus nicht
ableitbar.

Abbildung~\ref{fig:pa-panel-alurohr} zeigt das für den Testdatensatz erzeugte
Rohrmesh und die daraus abgeleitete lokale Centerline als Ergebnis der
Endstufen.

\begin{figure}[H]
  \centering
  \includegraphics[width=\textwidth]{pa_panel_alurohr_endprodukt.png}
  \caption{Alurohr-Endprodukt: rekonstruiertes Rohrmesh und die daraus
  extrahierte lokale Centerline. Beide Ansichten stammen aus demselben
  Entwicklungsstand (18. Juli 2026).}
  \label{fig:pa-panel-alurohr}
\end{figure}

Beim archivierten Referenzlauf extrahierte DGtal 87 Rohpunkte; der
Grad-10-B-Spline erzeugte 380 Ausgabepunkte. Diese Zahlen belegen die
Ausführbarkeit des Postprocessings. Sie bewerten weder die Lage der Linie im
Rohr noch ihre Formabweichung gegen eine Referenz.

\subsection{Negative Ergebnisse als Robustheitsbefund}

Mehrere zunächst fehlgeschlagene Läufe lieferten übertragbare Ergebnisse für
die Architektur. Die feste Eval-Liste muss Bildnamen unabhängig von
Dateiendungen abgleichen; die konkrete Eval-Maske muss nichtleer sein; ein
Unterprozess darf nicht die Konfigurations-TSV des Matrixrunners konsumieren;
und SuGaR-Renderings benötigen sowohl den richtigen Importpfad als auch das
erwartete Tensorlayout. Diese Korrekturen sind heute als Quality-Gates oder
explizite Fehlerprüfungen implementiert. Tabelle~\ref{tab:fehlerchronik}
dokumentiert weitere Fälle.
